% Part: second-order-logic
% Chapter: syntax-and-semantics
% Section: language-of-sol

\documentclass[../../../include/open-logic-section]{subfiles}

\begin{document}

\olfileid{sol}{syn}{lan}

\olsection{二阶逻辑的语言}

与一阶逻辑类似，二阶逻辑的表达式也是由包含\emph{变元}、
\emph{常项符号}、\emph{谓词符号}，有时还包括\emph{函数符号}的
基本词汇表构建而成的。由这些符号，连同逻辑联结词、量词以及
诸如括号和逗号的标点符号，便构成了\emph{项}和\emph{公式}。
不同之处在于，二阶逻辑除了包含关于对象的变元之外，还包含关于关系和函数的变元，
并且允许对这些变元进行量化。因此，二阶逻辑的逻辑符号包括一阶逻辑的全部逻辑符号，外加：

\begin{enumerate}
\item 对于每个元数~$n$，都有一个可数无穷的二阶关系变元集合：
  $\Obj V_0^n$, $\Obj V_1^n$, $\Obj V_2^n$, \dots
\item 一个可数无穷的二阶函数变元集合：
  $\Obj u_0^n$, $\Obj u_1^n$, $\Obj u_2^n$, \dots
\end{enumerate}

在一阶逻辑中，通常包含等词~$\eq$。同样在一阶逻辑中，语言~$\Lang{L}$ 的
非逻辑符号对于让我们能够表达任何有趣的内容至关重要。
当然，存在完全不使用非逻辑符号的语句，但如果只有~$\eq$，
就很难表达任何有趣的内容。
在二阶逻辑中，由于我们拥有不受限制的关系变元和函数变元供应，
因此即便没有专门的非逻辑符号，我们也能表达在一阶语言中所能表达的一切内容。

\end{document}